% BHU PYQ -- Manually transcribed & error-corrected
\documentclass[11pt,a4paper]{article}

\usepackage[a4paper,top=2.5cm,bottom=2.5cm,left=2.8cm,right=2.8cm,headheight=14pt]{geometry}
\usepackage{amsmath,amssymb}
\usepackage[T1]{fontenc}
\usepackage[utf8]{inputenc}
\usepackage{lmodern}
\usepackage{microtype}
\usepackage{fancyhdr}
\usepackage{enumitem}
\usepackage{hyperref}
\usepackage{lastpage}
\usepackage{parskip}

\hypersetup{colorlinks=true,urlcolor=black,linkcolor=black,
  pdftitle={MTB-504: Differential Geometry -- BHU PYQ},pdfauthor={Banaras Hindu University}}

\pagestyle{fancy}
\fancyhf{}
\renewcommand{\headrulewidth}{0.4pt}
\renewcommand{\footrulewidth}{0.4pt}
\fancyhead[L]{\small   \textit{Banaras Hindu University}}
\fancyhead[R]{\small   \textit{MTB-504: Differential Geometry}}
\fancyfoot[C]{\small Page  \thepage\ of \pageref{LastPage}}


\newlist{parts}{enumerate}{2}
\setlist[parts,1]{label=(\alph*),leftmargin=*,itemsep=5pt,topsep=3pt}
\setlist[parts,2]{label=(\roman*),leftmargin=*,itemsep=3pt,topsep=2pt}

\newcommand{\pts}[1]{\hfill   \textit{[#1]}}


\begin{document}


\begin{center}
  {\large\bfseries BANARAS HINDU UNIVERSITY}\\[2pt]
  {
\normalsize B.A./B.Sc. (Hons.) Semester V Examination, 2022-23}\\[6pt]
  \rule{0.8\linewidth}{0.5pt}\\[6pt]
  {\large\bfseries Paper No. MTB-504: Differential Geometry}\\[4pt]
  {
\normalsize Time: 3 Hours\hfill Maximum Marks: 70}
\end{center}

\rule{\linewidth}{0.4pt}

\smallskip



\noindent   \textit{Instructions: Answer   \textbf{five} questions including Question~1 which is compulsory. Symbols have their usual meanings.}

\medskip


\noindent   \textbf{Question 1.}

\begin{parts}
  \item Give the definition of a regular curve in Euclidean 3-space. Show that a curve $\alpha : \mathbb{R}  o \mathbb{R}^3$ defined by $\alpha(t) = (t^3, t^2, 0)$ for $t \in \mathbb{R}$ is not regular. \pts{2}
  \item Give the definition of normal plane and rectifying plane of a regular curve in $\mathbb{R}^3$. \pts{2}
  \item Give the definition of Frenet-Serret apparatus of a unit speed curve $\alpha : (a,b)  o \mathbb{R}^3$. \pts{2}
  \item State the local fundamental theorem for regular curves in Euclidean 3-space. \pts{2}
  \item Give the definitions of the first fundamental form and second fundamental form for a simple surface in Euclidean 3-space. \pts{2}
  \item Give the definitions of Christoffel symbols of the first kind and Christoffel symbols of the second kind for a simple surface in Euclidean 3-space. \pts{2}
  \item Give the definitions of Gaussian curvature and the mean curvature for a simple surface in Euclidean 3-space. \pts{2}
\end{parts}

\medskip


\noindent   \textbf{Question 2.}

\begin{parts}
  \item State and prove the Frenet-Serret formulas for a unit speed curve in Euclidean 3-space. \pts{7}
  \item Give the definition of a general helix in Euclidean 3-space. Prove that a unit speed curve $\alpha(s)$ with $\kappa(s) 
eq 0$ for every $s$ is a general helix if and only if $ \tau = c \kappa$ for some real constant $c$. \pts{7}
\end{parts}

\medskip


\noindent   \textbf{Question 3.}

\begin{parts}
  \item Give the definition of the osculating plane of a regular curve in Euclidean 3-space. For a unit speed curve $\alpha : (a,b)  o \mathbb{R}^3$, if $\kappa > 0$ and $ \tau 
eq 0$, then show that the equation of the osculating plane at a point $P$ is:
  \[ [R - \alpha, \alpha', \alpha''] = 0 \]
  where symbols have their usual meanings. \pts{7}
  \item Let $\alpha$ be a unit speed curve which lies on a sphere of radius $a > 0$ centered at the point $x_0$. If $ \tau 
eq 0$, then show that:
  \[ \alpha(s) - x_0 = -\rho \mathbf{n} - \rho' \sigma \mathbf{b} \]
  where $\rho = \frac{1}{\kappa}$ and $\sigma = \frac{1}{ \tau}$, and also show that:
  \[ \rho^2 + (\rho' \sigma)^2 = a^2 \]
  Do the normal planes of a spherical curve pass through the center? \pts{7}
\end{parts}

\medskip


\noindent   \textbf{Question 4.}

\begin{parts}
  \item State and prove the fundamental theorem of curves in Euclidean 3-space. \pts{7}
  \item Let $\alpha(s)$ be a unit speed curve with $\kappa 
eq 0$. Show that the following three statements are equivalent:
  
\begin{parts}
    \item $\alpha(s)$ is a planar curve,
    \item the binormal vector $\mathbf{b}(s)$ is a constant vector, and
    \item the torsion $ \tau(s) = 0$ for all $s$.
  \end{parts} \pts{7}
\end{parts}

\medskip


\noindent   \textbf{Question 5.}

\begin{parts}
  \item Give the definition of a simple surface of class $C^k$. Prove that $\mathbf{x} : \mathbb{R}^2  o \mathbb{R}^3$ given by
  \[ \mathbf{x}(u,v) = \left( \frac{2au}{1+u^2+v^2}, \frac{2av}{1+u^2+v^2}, \frac{u^2+v^2-1}{1+u^2+v^2} a \right) \]
  is a simple surface which represents a part of a sphere. \pts{7}
  \item Give the definition of the tangent plane at a point of a simple surface. Let
  \[ \mathbf{x}(u,v) = \left( u, v, \sqrt{1-u^2-v^2} \right) \]
  be a simple surface. Obtain the tangent plane at the point $\mathbf{x}(1/2, 1/2)$. \pts{7}
\end{parts}

\medskip


\noindent   \textbf{Question 6.}

\begin{parts}
  \item Give the definition of the area $S(D)$ of a region $D$ covered by a surface patch $\mathbf{x} : U  o \mathbb{R}^3$. Obtain the area of a sphere of radius $a > 0$, given by
  \[ \mathbf{x}(u,v) = (a \sin u \cos v, a \sin u \sin v, a \cos u) \] \pts{7}
  \item Obtain the Christoffel symbols of the second kind for the simple surface $\mathbf{x}(u, v) = (u, v, uv)$. \pts{7}
\end{parts}

\medskip


\noindent   \textbf{Question 7.}

\begin{parts}
  \item State and prove the Gauss and Weingarten formulas for a simple surface $\mathbf{x} : U  o \mathbb{R}^3$ of class $\ge C^2$. \pts{7}
  \item Obtain the Gaussian curvature of a sphere of radius $a > 0$ in Euclidean 3-space. \pts{7}
\end{parts}




\vfill
\rule{\linewidth}{0.4pt}

\begin{center}\small
\href{https://bhu-exam.vercel.app/}{Click here to visit our website}\\[4pt]\href{https://me-aryan.vercel.app/}{Click here to visit our contributor}\\[4pt]\href{https://quashan-me.vercel.app/}{Click here to visit our contributor}
\end{center}

\end{document}
